% Part: methods
% Chapter: proofs
% Section: cant-do-it

\documentclass[../../../include/open-logic-section]{subfiles}

\begin{document}

\olfileid[id]{mth}{prf}{cnt}

\olsection{Saya Tidak Bisa Melakukannya!{}}

Kita semua pernah sampai pada titik ketika rasanya ingin menyerah. Namun,
Anda \emph{bisa} melakukannya. Dosen dan asisten pengajar Anda, serta
sesama mahasiswa, dapat membantu. Mintalah bantuan mereka!{} Berikut
beberapa kiat untuk membantu Anda menghindari krisis, serta tindakan yang
dapat diambil jika Anda merasa ingin menyerah.

Untuk memastikan bahwa Anda dapat menyelesaikan soal dengan baik, lakukan
hal-hal berikut:
\begin{enumerate}
\item \emph{Mulailah sedini mungkin.} Kita menjadi sibuk sepanjang semester,
  dan banyak di antara kita bergumul dengan kebiasaan menunda-nunda. Salah
  satu hal terbaik yang dapat Anda lakukan ialah mulai mengerjakan tugas
  sejak dini. Dengan demikian, jika menemui jalan buntu, Anda masih punya
  waktu untuk mencari penyelesaian (yang bukan menangis).
\item \emph{Berbicaralah dengan teman sekelas Anda.} Anda tidak sendirian.
  Orang lain di kelas mungkin juga mengalami kesulitan---tetapi kesulitan
  mereka mungkin berbeda. Membicarakannya dengan rekan-rekan dapat memberi
  Anda sudut pandang lain terhadap soal, yang mungkin menghasilkan terobosan.
  Tentu saja, jangan sekadar menyalin penyelesaian mereka: mintalah petunjuk,
  atau jelaskan letak kebuntuan Anda dan tanyakan langkah berikutnya. Ketika
  Anda akhirnya memahaminya, balaslah bantuan itu. Membantu orang lain
  melangkah maju dan menjelaskan berbagai hal juga akan membantu Anda
  memahaminya dengan lebih baik.
\item \emph{Mintalah bantuan.} Ada banyak sumber bantuan yang tersedia bagi
  Anda---dosen dan asisten pengajar siap membantu dan \emph{ingin} Anda
  berhasil. Mereka semestinya dapat membantu Anda menguraikan suatu soal dan
  mengenali bagian proses yang membuat Anda kesulitan.
\item \emph{Beristirahatlah.} Jika menemui jalan buntu, hal itu \emph{mungkin}
  terjadi karena Anda sudah terlalu lama menatap soal tersebut. Beristirahatlah
  sebentar, minumlah secangkir teh, atau kerjakan soal lain untuk sementara,
  kemudian kembalilah ke soal itu dengan pikiran segar. Tidurlah dahulu dan
  pikirkan lagi sesudahnya.
\end{enumerate}

Perhatikan bahwa strategi-strategi ini mengharuskan Anda sudah mulai
mengerjakan bukti jauh-jauh hari. Jika Anda baru mulai mengerjakan bukti
pukul dua dini hari pada hari sebelum tenggatnya, strategi-strategi ini
mungkin tidak terlalu membantu.

Semua ini mungkin terdengar muram, tetapi menyelesaikan sebuah bukti adalah
tantangan yang pada akhirnya memberikan kepuasan. Ada orang yang menjadikan
hal ini sebagai profesi---jadi tentu ada sesuatu yang dapat dinikmati di
dalamnya. Seperti hampir segala hal lain, menyelesaikan soal dan menyusun
bukti memerlukan latihan. Anda mungkin melihat teman sekelas yang tampaknya
mudah melakukannya: barangkali mereka memang sudah banyak berlatih. Usahakan
agar tidak terlalu mudah menyerah.

Jika Anda memang kehabisan waktu (atau kesabaran) untuk suatu soal tertentu,
tidak apa-apa. Itu tidak berarti bahwa Anda bodoh atau tidak akan pernah
memahaminya. Cari tahu (dari dosen atau mahasiswa lain) bagaimana soal itu
diselesaikan, lalu kenali letak kesalahan atau kebuntuan Anda agar hal yang
sama dapat dihindari ketika kelak menghadapi persoalan serupa. Setelah itu,
cobalah mengerjakannya tanpa melihat penyelesaian. Dan lain kali, mulailah
(serta mintalah bantuan) lebih awal.

\end{document}
